%%%%%%%%%%%%%%%%%%%%%%%%%%%%%%%%%%%%%%%%%%%%%%%%%%%%%%%%%%%%%%
%%%%%%%%%%%%%%%%%% MA-0525 Combinatoria %%%%%%%%%%%%%%%%%%%%%%%%%%%%%%
%%%%%%%%%%%%%%%%%%%%%%%%%%%%%%%%%%%%%%%%%%%%%%%%%%%%%%%%%%%%%%
\newpage
\subsection*{MA-0525 Combinatoria}
\addcontentsline{toc}{subsection}{MA-0525 Combinatoria}

\hypertarget{MA0525}{}
\begin{refsection}
\begin{table}[H]
\centering{
\begin{tabular}{|ll|}
\hline
\textbf{Año:} IV  & \textbf{Requisitos:} MA-0351 Cálculo en varias variables \\
& y MA-0361 Álgebra Lineal I \\ 
\textbf{Ciclo:} Optativa  & \textbf{Correquisitos:} Ninguno \\ 
\textbf{Tipo de curso:} Teórico & \textbf{Horas presenciales por semana:} 5 \\ 
\textbf{Créditos:} 4 & \textbf{Horas de trabajo independiente por semana:} 10 \\
\textbf{Virtualidad:} P  &  \\ \hline
\end{tabular}
}
\end{table}

\subsubsection*{Descripción del curso}

Un tema común de la combinatoria es que estudia principalmente estructuras finitas o discretas; sin embargo, la mejor forma de describir esta área es por medio de los problemas que aborda.

En combinatoria, las técnicas y métodos para resolver un problema suelen ser más importantes que el resultado mismo. En este curso se estudian herramientas básicas de conteo y nociones fundamentales de matemática discreta, como los coeficientes binomiales, los grafos y las funciones generatrices. Además, se introducen problemas de combinatoria extremal.

\subsubsection*{Objetivos}

\textbf{General:} Desarrollar en la persona estudiante competencias para abordar problemas combinatorios mediante técnicas de conteo, teoría de grafos, métodos probabilísticos y combinatoria extremal. \\

\textbf{Específicos:} Al finalizar el curso se espera que la persona estudiante sea capaz de:

\begin{enumerate}
\item Aplicar herramientas y resultados fundamentales de la combinatoria enumerativa, desde los primeros principios de conteo hasta el uso de funciones generatrices.
\item Utilizar nociones básicas de la teoría de grafos y sus herramientas en problemas combinatorios.
\item Aplicar ideas de combinatoria extremal y del método probabilístico en la demostración de resultados y la resolución de problemas.
\item Seleccionar y combinar estrategias y técnicas habituales en la resolución de problemas combinatorios.
\end{enumerate}

\subsubsection*{Contenidos}

\begin{enumerate}
\item \textbf{Herramientas de conteo (3 semanas):}
\begin{enumerate}
    \item Regla de la suma y del producto; conteo de palabras; acomodos con y sin repetición; permutaciones.
    \item Coeficientes binomiales; teorema del binomio de Newton; identidades con coeficientes binomiales.
    \item Principio de inclusión--exclusión.
    \item Énfasis en biyecciones para transformar un conteo en otro.
\end{enumerate}

\item \textbf{Introducción a la teoría de grafos (2{,}5 semanas):}
\begin{enumerate}
    \item Nociones básicas de grafos; lema del apretón de manos.
    \item Matrices de adyacencia; estructuras de incidencia.
    \item Árboles; emparejamientos; número cromático.
\end{enumerate}

\item \textbf{Método probabilístico (2 semanas):}
\begin{enumerate}
    \item Principio de palomar y generalizaciones.
    \item Espacio de probabilidad; variables aleatorias discretas; esperanza; desigualdad de Markov.
    \item Aplicaciones del método probabilístico.
    \item Énfasis en probabilidad discreta, con ejemplos geométricos elementales cuando sea pertinente.
\end{enumerate}

\item \textbf{Funciones generatrices (2{,}5 semanas):}
\begin{enumerate}
    \item Propiedades y series conocidas; conteos con funciones generatrices.
    \item Particiones; relaciones de recurrencia.
    \item Función generatriz exponencial; identidades con funciones generatrices.
\end{enumerate}

\item \textbf{Combinatoria extremal (2{,}5 semanas):}
\begin{enumerate}
    \item Introducción a problemas extremales.
    \item Problemas extremales en conjuntos: teorema de Sperner; teorema de Erd\H{o}s--Ko--Rado.
    \item Problemas extremales en grafos: teorema de Mantel; teorema de Turán; teorema de Ramsey.
\end{enumerate}
\end{enumerate}

\subsubsection*{Temas adicionales opcionales}

Según el ritmo del grupo y el interés del cuerpo docente, pueden desarrollarse módulos complementarios, por ejemplo:
\begin{enumerate}
\item \textbf{Teoría extremal de grafos:} número extremal de un grafo; número extremal de árboles y ciclos; caracterización a partir del número cromático; teoremas de estabilidad y supersaturación.
\item \textbf{Teoría de Ramsey:} teorema de van der Waerden; teorema de Schur; teorema de Erd\H{o}s--Szekeres.
\item \textbf{Aplicaciones de álgebra lineal en combinatoria:} conjuntos con intersecciones pares; desigualdad de Fisher; desigualdades de Frankl--Wilson y de Bollobás.
\item \textbf{Relaciones entre topología y combinatoria:} lema de Tucker; teorema de Borsuk--Ulam; grafos de Kneser; lema de Sperner y punto fijo de Brouwer.
\item \textbf{Método polinomial:} Combinatorial Nullstellensatz; conjuntos de Nikodym y Kakeya en cuerpos finitos.
\end{enumerate}

\subsubsection*{Metodología}

\noindent\textbf{Lineamientos metodológicos:} La persona docente a cargo de este curso debe respetar las siguientes pautas:
\begin{enumerate}
    \item Combinar exposición teórica con resolución guiada y autónoma de problemas combinatorios.
    \item Priorizar el razonamiento combinatorial (biyecciones, doble conteo, argumentos extremales) por encima de la memorización de fórmulas.
    \item Promover la comunicación clara de demostraciones y estrategias, en forma oral y escrita.
    \item Fomentar el trabajo independiente mediante listas de ejercicios y lecturas selectas de la bibliografía.
\end{enumerate}

\textbf{Sugerencias metodológicas:} Adicionalmente, se recomienda:
\begin{enumerate}
    \item Desarrollar talleres de problemas que integren conteo, grafos, probabilidad y funciones generatrices.
    \item Utilizar software o entornos simbólicos para explorar ejemplos de grafos, conteos y funciones generatrices cuando apoye la comprensión conceptual.
    \item Relacionar los temas opcionales con intereses del grupo (ciencia de datos, optimización, teoría de números, etc.) cuando sea pertinente.
\end{enumerate}

\nocite{AloSpe16}
\nocite{And02}
\nocite{AndFen04}
\nocite{BabFra20}
\nocite{Bol86}
\nocite{Bol02}
\nocite{BonMur07}
\nocite{Bru04}
\nocite{FueSim13}
\nocite{GraRotSpe90}
\nocite{Juk11}
\nocite{KelTro17}
\nocite{Mat03}
\nocite{MorOli11}
\nocite{Sob13}
\nocite{TaoVu06}
\nocite{Vil72}
\nocite{Zam25}

\printcursosbibliography

\end{refsection}
